\section{General algorithmic framework for approximating \textsc{HC-$\cF$}}

Having characterized the kinds of graph classes that will fit into our framework we can now move on towards desciribing how to use the properties of $\alpha\br{n}$-well-behavedness to construct a general approximation algorithm for \textsc{HC-$\cF$}. The main idea is that we use a level-based \textsc{LP} relaxation to guide the recursive algorithm.
In the preprocessing step we encode the problem using an \textsc{ILP}. We obtain this formulation by combining any $\alpha\br{n}$-well-behaved \textsc{ILP} with the level decomposition approach of Charikar et. al in \cite{ApproximateHierarchicalClusteringViaSparsestCutAndSpreadingMetrics}. Our \textsc{ILP} consists of multiple variable levels which intuitively encode levels of the clustering tree.

After solving \textsc{LP-HC-$\cF$}, we use a rounding routine to get a partially integral solution which at each level corresponds to choosing which vertices should satisfy the $\cF$-predicate.

During the recursive step, we process clusters in a top-down manner starting from $H=G$. For a cluster $H$, we set $t=w\br{V\br{H}}/2$. Then, we look which vertices satisfy the $\cF$-predicate at level $t$ and which do not.
\DDfix{(For a formal level definition see Section~\ref{sec:levels} below).}
According to this, we partition them into two sets $X_1$ (satisfying the predicate) and $X_0$ (not satisfying the predicate). The solution of \textsc{LP-HC-$\cF$} at level $t$ is then used as a valid solution to \textsc{$p_{\cF}$-P-LP} with $X=X_1$, yielding a cut $F_H$ with $\alpha\br{n}$ loss. The same level-$t$ fragment is also used as a valid guide for \textsc{$\rho$-SEP} with $\rho=1/2$, $\rho_0=2/3$ and $X=X_0$, yielding a cut $B_H$ such that each connected component of $H\setminus B_H$ that contains a terminal has total weight at most $2w\br{V\br{H}}/3$ and the cut cost is within a factor $\beta\br{n}=O\br{\log n}$ of the guide. After cutting edges in $F_H\cup B_H$, we recursively process all connected components.

Since we decompose the \textsc{LP} into separate levels, by summing up over all of them we will be able to show that our algorithm achieves an approximation ratio of $O\br{\alpha\br{n}+\beta\br{n}}$, where $\beta\br{n}$ is the approximation ratio of the \textsc{$\rho$-SEP} rounding algorithm.

\subsection{Level decomposition of \textsc{HC-$\cF$}} \label{sec:levels}

In order to analyze the structure of \textsc{HC-$\cF$} trees, we introduce the notion of level decomposition. For a given \textsc{HC-$\cF$} tree $T$, we define the \emph{level decomposition} as follows.
Let $\cC^T$ denote the family of all clusters in $T$, $\cC^T=\brc{G_u\colon u\in V(T)}$.
We assume that vertex weights are integral, hence $w\br{G}\in \mathbb{N}^+$. For $t\in [w\br{G}]$, denote by $\cL_t^T$ the set of all maximal clusters in $\cC^T$ of total weight at most $t$. We call $\cL_t^T$ the $t$-th \emph{level} of the decomposition.
Observe that for each $t$, $\cL_t^T$ is a partition of $V\br{G}$. By $S_t^T$ we denote the set of edges cut at level $t+1$, that is the set of edges with endpoints in different clusters of $\cL_t^T$. The shift of indices is for the sake of convenience. We call $S_t^T$ the $t$-th \emph{level cut}.
It is easy to see that the cost of \textsc{HC-$\cF$} tree $T$ can be expressed as follows:
\begin{observation}\label{obs:cost-decomposition}
    $\COST_G\br{T} = \sum_{t=1}^{w\br{G}} c\br{S_t^T}$.
\end{observation}

Since we use the above levels to decompose the clustering tree and its cost, we would also like to do the same for the $\cF$-predicate. This is because eventually, all vertices satisfy the $\cF$-predicate, but whether a vertex $v\in V$ satisfies the $\cF$-predicate, when it belongs to a cluster of weight $t$ depends on the structure of the clustering tree.
For every level $t\in [w\br{G}]$ and every vertex $v\in V\br{G}$, we define the level-indexed $\cF$-predicate indicator by setting $p_{\cF}^t\br{v}$ to be true iff $p_{\cF}\br{v}$ is true in $\cL_t^T$. 
\subsection{\textsc{ILP} formulation of \textsc{HC-$\cF$} and its relaxation}
    We nowe move to a description of the constraints of the \textsc{ILP} along with an explanation of the role they play. We use the well-known \emph{spreading-metrics} framework which in recent years have come particularly useful for approximating many problems regarding graph cutting \DD{[citation here?] (aby nie zapomniec)}. The idea is that we encode \textsc{HC-$\cF$} as a level-indexed formulation with edge-length variables and shortest-path distances.
    For each level $t$, each cluster is either of total weight at most $t$ or it is an $\cF$-cluster. The partition at level $t-1$ refines the partition at level $t$.
    For each edge $e\in E$ and level $t\in [w\br{G}]$ we use a variable $x_e^t\in[0,1]$. Let $\dist_{x}^{t}(\cdot,\cdot)$ denote shortest-path distances in $G$ induced by lengths $\{x_e^t\}_{e\in E}$.
    \DD{czyli musimy tutaj wykladniczo enumerowac wszystkie sciezki brac min?}
    We also define a variable $y_v^t=1$ iff $p_{\cF}^t\br{v}$ is true and $0$ otherwise.
    \DD{tutaj nalezaloby skomentowac, ze wczesniej dopuszczamy tylko edge cut variables, a tutaj jednak mamy jeszcze inne.}
    We first explain all constraints and then we give the full ILP. 
    These are as follows:
    \begin{enumerate}
        \item \textbf{refinement:} $x_{e}^t \leq x_{e}^{t-1}$ for all $e\in E$ and $t\in [w\br{G}]$, because the partition at level $t-1$ refines the partition at level $t$. Moreover, $y_v^t \leq y_v^{t-1}$ for all $v\in V$ and $t\in [w\br{G}]$, since $p_{\cF}^t\br{v}$ implies $p_{\cF}^{t-1}\br{v}$.

        \item \textbf{$p_{\cF}$-P constraints:} Recall, that for every vertex $v\in V$, the set $\cR_v$ contains all subsets of edges $R\subseteq E$ such that if at least one edge in $R$ is cut, then $p_{\cF}\br{v}$ is true. Therefore, for each $v\in V$, every level $t\in [w\br{G}]$ and every $R\in \cR_v$ we have $\sum_{e\in R} x_e^t \geq y_{v}^t$. Additionally, we have $y_v^1=1$, since $p_{\cF}^1\br{v}$ is true for all $v\in V$.

        \item \textbf{spreading-metrics:} For every vertex $v\in V$, every level $t\in [w\br{G}]$, and every subset $S\subseteq V$ containing $v$, we have
        \[
            \sum_{u\in S} w(u)\cdot \dist_{x}^{t}(v,u) \geq \br{1-y_v^t}\cdot \br{w\br{S}-t}.
        \]
        When $y_v^t=0$, this condition enforces that inside any set $S$ containing $v$, enough weighted mass is separated from $v$ unless the cluster weight around $v$ is already at most $t$. When $y_v^t=1$, the right-hand side vanishes, corresponding to the fact that the cluster containing $v$ is an $\cF$-cluster and hence its weight can be larger than $t$.
    \end{enumerate}

    The following \textsc{ILP}, which we call \textsc{ILP-HC-$\cF$}, is an exact formulation of the \textsc{HC-$\cF$} problem:
    \begin{align}
        \min \quad &\sum_{t=1}^{w\br{G}} \sum_{e\in E} c(e) \cdot x_{e}^t\\
        \text{subject to} \quad &x_{e}^t \leq x_{e}^{t-1} \quad \forall e\in E,\; t\in [w\br{G}],\\
        &y_v^t \leq y_v^{t-1} \quad \forall v\in V, t\in [w\br{G}],\\
        &y_v^1 = 1 \quad \forall v\in V,\\
        &\sum_{e\in R} x_{e}^t \geq y_{v}^t \quad \forall v\in V, R\in \cR_v, t\in [w\br{G}], \label{cover_constraint_ILP}\\
        &\sum_{u\in S} w(u)\cdot \dist_{x}^{t}(v,u) \geq \br{1-y_v^t}\cdot \br{w\br{S}-t} \quad \forall v\in V,\; S\subseteq V,\; v\in S,\; t\in [w\br{G}], \label{spreading_constraint_ILP}\\
        &x_{e}^t \in \{0,1\} \quad \forall e\in E, t\in [w\br{G}],\\
        &y_v^t \in \{0,1\} \quad \forall v\in V, t\in [w\br{G}].
    \end{align}

    \begin{lemma}\label{lem:ilp-exact}
        \textsc{ILP-HC-$\cF$} is an exact formulation of \textsc{HC-$\cF$}. In particular, the optimum \textsc{ILP} value equals the optimum hierarchical-clustering value.
        \DD{moze dobrze byloby wprowadzic symbol na optimum naszego problemu, np. $\OPT_{\cF}(G)$; wowczas aby bylo spojnie dotychczasowe $\OPT_G(T)$ staloby sie $\OPT_{\cF}(T)$, i mielibysmy $\OPT_{\cF}(G):=\min\brc{\OPT_{\cF}(T)\colon ...}$.}
    \end{lemma}
    \begin{proof}
        $(\Rightarrow)$ Let $T$ be any feasible \textsc{HC-$\cF$} tree. \DD{moze tez dobrze miec formalna fraze na te drzewa, np. $\cF$-clustering tree zdefiniowac?} For each level $t\in [w\br{G}]$, let $S_t^T$ be the level cut from the level decomposition and set
        \[
            x_e^t:=\mathbf{1}[e\in S_t^T].
        \]
        Also set $y_v^t:=1$ iff the level-$t$ cluster of $v$ satisfies $p_{\cF}$.

        The refinement constraints hold because level cuts are nested. The constraints on $y$ hold by monotonicity of the predicate along the refinement. The cover constraints hold by definition of the family $\cR_v$: whenever $y_v^t=1$, each required cover-type constraint for $v$ at level $t$ is satisfied by the cut pattern encoded by $x^t$.

        For spreading constraints, fix $v\in V$, $t\in [w\br{G}]$, and $S\subseteq V$ with $v\in S$. If $y_v^t=1$, the right-hand side is $0$. Assume $y_v^t=0$ and let $C$ be the level-$t$ cluster containing $v$. Then $C$ is not an $\cF$-cluster, so by level definition $w(C)\le t$. Every $u\in S\setminus C$ is separated from $v$ by level-$t$ cuts, hence $\dist_x^{t}(v,u)\ge 1$. Therefore,
        \[
            \sum_{u\in S} w(u)\cdot \dist_x^{t}(v,u)
            \ge \sum_{u\in S\setminus C} w(u)
            = w(S)-w(S\cap C)
            \ge w(S)-w(C)
            \ge w(S)-t,
        \]
        which is exactly \eqref{spreading_constraint_ILP} for $y_v^t=0$. Thus $(x,y)$ is ILP-feasible and
        \[
            \sum_{t=1}^{w\br{G}}\sum_{e\in E} c(e)x_e^t
            = \sum_{t=1}^{w\br{G}} c(S_t^T)
            = \COST_G(T)
        \]
        by Observation~\ref{obs:cost-decomposition}. Hence $\OPT_{ILP}\le \OPT_{HC}$.
        \DD{widze, ze tutaj OPT-y sie pojawily. Proponuje zdecydowac sie na jakis wariant, dac def. do preliminaries lub wczesniej.}

        $(\Leftarrow)$ Let $(x,y)$ be any integral feasible ILP solution. For each level $t$, define
        \[
            E_t^0:=\brc{e\in E\colon x_e^t=0},
        \]
        and let $\Pi_t$ be the partition of \DDfix{$V$ into connected components of $G\setminus \br{E\setminus E_t^0}$. (Note that $E\setminus E_t^0$ are exactly the edges cut at level $t$).} Because $x_e^t\le x_e^{t-1}$, we have $E_{t-1}^0\subseteq E_t^0$, so $\Pi_{t-1}$ refines $\Pi_t$ and the family $(\Pi_t)_{t=1}^{w\br{G}}$ is laminar; therefore it defines a hierarchical tree $T$.
        \DD{tak sie zastanawiam czy nie powinnismy miec warunkow brzegowych na $x$-y (np. $x_e^t=0$ for $t=\spr{w(G)}$? Bo wtedy tutaj mamy wlasnosc, ze $E_t^0=E$ dla $t=\spr{w(G)}$ i wtedy mamy root corresponds to $G$.}

        Consider any level $t$ and any component $C\in \Pi_t$. For every $v\in C$, all vertices of $C$ are at zero $\dist_x^t$-distance from $v$, so plugging $S=C$ into \eqref{spreading_constraint_ILP} gives
        \[
            0\ge (1-y_v^t)\cdot (w(C)-t).
        \]
        Hence either $w(C)\le t$, or $y_v^t=1$ for all $v\in C$. In the second case, by cover constraints and the definition of $\cR_v$, each vertex in $C$ satisfies the $\cF$-predicate; therefore $C\in \cF$. So every cluster in level $t$ is either of weight at most $t$ or belongs to $\cF$, exactly as required \DDfix{in a clustering tree.}

        Finally, the level cut induced by $\Pi_t$ is precisely $\brc{e\in E\colon x_e^t=1}$, so
        \[
            \COST_G(T)=\sum_{t=1}^{w\br{G}} c\br{\brc{e\in E\colon x_e^t=1}}
            =\sum_{t=1}^{w\br{G}}\sum_{e\in E} c(e)x_e^t.
        \]
        Therefore every integral ILP solution maps to a feasible \textsc{HC-$\cF$} tree of the same cost, implying $\OPT_{HC}\le \OPT_{ILP}$, which completes the proof.
        \DD{terminologia}
    \end{proof}

    Observe, that the \textsc{LP} relaxation \textsc{LP-HC-$\cF$} of the above \textsc{ILP} may have an exponential number of constraints, due to the constraints of type \eqref{cover_constraint_ILP}, however we are guaranteed that a separation oracle exists for these constraints. Therefore, we can solve this \textsc{LP} in polynomial time using the ellipsoid method.
    %Note that in our relaxation we allow each $x_e^t\geq 0$ and each $y_v^t$ to take any value in $[0,1]$.
    Let $LP$ be a solution to \textsc{LP-HC-$\cF$} and let $c\br{LP}$ be the value of its objective function. We now show how to round the solution of \textsc{LP-HC-$\cF$} to get a partially integral solution without excessive increase in cost.

    \begin{lemma}\label{lem:lp-rounding}
        Round the values of $y_v^t$ to get $\hat{y}_v^t = 1$ if $y_v^t \geq 1/2$ and $\hat{y}_v^t = 0$ otherwise
        and $\hat{x}_{e}^t=2\cdot x_{e}^t$. Then, $\hat{x}, \hat{y}$ is a feasible solution for \textsc{LP-HC-$\cF$} with cost at most $2\cdot c\br{LP}$.
    \end{lemma}
    \begin{proof}
        Firstly, observe that the refinement constraints are satisfied since $\hat{x}_{e}^t \leq \hat{x}_{e}^{t-1}$ and $\hat{y}_v^t \leq \hat{y}_v^{t-1}$ for all $e\in E$, $v\in V$ and $t\in [w\br{G}]$. The cover \DD{(wczesniej nazwalismy je $p_{\cF}$-P constraints, wiec markuje do ujednolicenia)} constraints are satisfied since if $\hat{y}_v^t=1$, then $y_v^t\geq 1/2$ and hence $\sum_{e\in R} x_{e}^t \geq y_v^t \geq 1/2$ for all $R\in \cR_v$. Therefore, $\sum_{e\in R} \hat{x}_{e}^t = 2\cdot \sum_{e\in R} x_{e}^t \geq 1$. If $\hat{y}_v^t=0$, then the cover constraint is trivially satisfied. The spreading-metrics constraints are satisfied since if $\hat{y}_v^t=0$, then $y_v^t<1/2$ and hence
        \[
            \sum_{u\in S} w(u)\cdot \dist_{x}^{t}(v,u) \geq \br{1-y_v^t}\cdot \br{w\br{S}-t} > 1/2\cdot \br{w\br{S}-t}.
        \]
        Therefore, we get that for all $S\subseteq V$ containing $v$, such that $\bar{y}_v^t=0$, we have
        \[
            \sum_{u\in S} w(u)\cdot \dist_{\hat{x}}^{t}(v,u) = 2\cdot \sum_{u\in S} w(u)\cdot \dist_{x}^{t}(v,u) > w\br{S}-t
        \]
        as required.
        If $\hat{y}_v^t=1$, then the spreading-metrics constraint is trivially satisfied.
        Finally, the cost of the new solution is at most $2\cdot c\br{LP}$ since $\sum_{t=1}^{w\br{G}} \sum_{e\in E} c(e) \cdot \hat{x}_{e}^t = 2\cdot \sum_{t=1}^{w\br{G}} \sum_{e\in E} c(e) \cdot x_{e}^t = 2\cdot c\br{LP}$.
    \end{proof}

    Overall, we obtain a partially integral point where all of the $y$ values are integral, which translates to the fact that we have decided which vertices need to satisfy the $\cF$-predicate at each level of the clustering tree. From now on we will thus assume that for every $v\in V$ and $t\in[w\br{G}]$, the value of $p_{\cF}^t\br{v}$ is fixed and we can remove all of the $y$ variables from the \textsc{LP}. Since we will not further use the original values, by a slight abuse of notation from now on we will refer to this solution simply by $LP$.
\subsection{Restricting the solutions}
    
In our charging scheme, \DD{fraza charging scheme pojawia sie pierwszy raz} we will need to split the $LP$ into several solutions restricted to the clusters of the hierarchical clustering tree $T$ we are building. Fix some cluster $H$ of $G$ and a level $t$. By $LP_t^H$ denote the solution of \textsc{LP-HC-$\cF$} restricted to the set of vertices $H$ at level $t$ which is done by discarding all of the variables $x_{uv}^t$ concerning vertices outside of $H$ or levels other than $t$. We call each such restricted part $LP_t^H$ a \emph{guide}. For any target problem $\Pi$ with LP relaxation \textsc{LP-$\Pi$} (on the same graph fragment and terminal set), we say that $LP_t^H$ is a \emph{valid guide} for $\Pi$ if the restricted values from $LP_t^H$ satisfy all constraints of \textsc{LP-$\Pi$}. In other words, a valid guide is simply a feasible LP point for the target problem, inherited from \textsc{LP-HC-$\cF$}. This restriction is also performed on the cost function and this is reflected in the accordingly defined quantities $\COST\br{LP_t^H} = \sum_{uv\in E\br{H}} c\br{u,v} \cdot x_{uv}^t$. We have the following lemma:

    \begin{lemma}\label{lem:relaxation-decomposition}
        Let $T$ be any hierarchical clustering tree. For any cluster $H \in \cC^T$, let $r\br{H}$ be the weight of the heaviest child cluster of $H$ if it exists or $0$ otherwise. Then:
        \begin{align*}
            \sum_{H\in \cC^T}\sum_{t=r\br{H}+1}^{w\br{V\br{H}}} c\br{LP_t^H}\leq c\br{LP}
        \end{align*}
    \end{lemma}
    \begin{proof}
        It suffices to show that for any fixed edge $uv\in E$ and level $t\in [w\br{G}]$, the term $c\br{u,v}\cdot x_{uv}^t$ is counted at most once in the above sum. Assume for contradiction that there are two clusters $H$ and $B$ both contributing this term at level $t$. By definition of hierarchical clustering, one of $H$, $B$ must be an ancestor of the other; without loss of generality assume $H$ is an ancestor of $B$. Then
        \begin{align*}
            t &\in \left(r\br{H}, w\br{V\br{H}}\right],\\
            t &\in \left(r\br{B}, w\br{V\br{B}}\right].
        \end{align*}
        Since $H$ is an ancestor of $B$, we have
        \begin{align*}
            r\br{H} &\geq w\br{V\br{B}}.
        \end{align*}
        Hence the above intervals are disjoint, a contradiction. The bound follows by summing over all edges and levels.
    \end{proof}

    \begin{lemma}\label{lem:lp-level-bound}
        Let $H$ be any cluster and let $d\in\mathbb{N}^+$. Then
        \[
        w\br{V\br{H}}\cdot c\br{LP_{\br{1-1/d}w\br{V\br{H}}}^H}\leq d\cdot \sum_{t=\br{1-1/d}w\br{V\br{H}}+1}^{w\br{V\br{H}}} c\br{LP_t^H}.
        \]
    \end{lemma}
    \begin{proof}
    The inequality follows from monotonicity of level cuts: as $t$ increases, more edges belong to the cut, so averaging over the last $w\br{V\br{H}}/d$ levels gives the claim.
    \end{proof}

    Fix a subgraph $H$ of $H$ with weight $r=w\br{V\br{H}}$ and a level $t\in [r,w\br{V\br{H}}]$. Define
    \[
        X_1:=\brc{v\in V\br{H}\colon \hat{y}_v^t=1}, \qquad X_0:=\brc{v\in V\br{H}\colon \hat{y}_v^t=0}.
    \]

    The next lemma allows us to decompose the solution to \textsc{LP-HC-$\cF$} into several instances of the \textsc{LP-$p_{\cF}$-P}:
    \begin{lemma}\label{lem:lp-fes}
        Let $H$ be a non $\cF$-cluster of weight $r=w\br{V\br{H}}$ and let $r\leq t$ be some fixed level. Then $LP_r^H$ is a valid solution for the \textsc{$p_{\cF}$-P} for $X_1$.
    \end{lemma}
    \begin{proof}
        In $LP_r^H$ there are additional metric/spreading constraints for vertices $v$ such that $p_{\cF}^t(v)=0$. This is acceptable since additional constraints yield a stricter relaxation. Now consider any $v\in V$. If $p_{\cF}^t(v)=0$, then the cover constraint is always satisfied, which corresponds to the fact that $v$ does not need to belong to a $\cF$-cluster. If $p_{\cF}^t(v)=1$, then the cover constraint is satisfied since it is explicitly included in the \textsc{LP-HC-$\cF$} formulation. Moreover, since $v\in X_1$, we have that all cover constraints in the \textsc{LP-$p_{\cF}$-P} formulation are included in the \textsc{LP-HC-$\cF$} formulation. Therefore, the claim follows.
    \end{proof}

    
    We also need a decomposition statement for the size-reduction primitive. The next lemma states that the same level-restricted \textsc{LP-HC-$\cF$} solution can be viewed as a feasible guide for \textsc{$\rho$-SEP}.
    \begin{lemma}\label{lem:lp-ksep}
        Let $H$ be a non $\cF$-cluster of weight $r=w\br{V\br{H}}$ and let $t=r/2$. Then $LP_t^H$ is a valid guide for the weighted spreading-metrics \textsc{LP-$\rho$-SEP} with $\rho=1/2$ for $X_0$.
    \end{lemma}
    \begin{proof}
        The statement is analogous to Lemma~\ref{lem:lp-fes}, but with complement terminal set.
        First, the cover-type constraints are irrelevant for \textsc{$\rho$-SEP} feasibility and can only make the relaxation stronger.

        Second, for vertices in $X_0$ (i.e., with $p_{\cF}^t(v)=0$), the formulation contains exactly the subset-based spreading constraints needed for \textsc{LP-$\rho$-SEP}: for every $S\subseteq V(H)$ and every $v\in S\cap X_0$,
        \[
            \sum_{u\in S} w(u)\cdot \dist_{x}^{t}(v,u) \geq w\br{S}-t.
        \]
        With $t=r/2$, this is exactly the weighted spreading requirement corresponding to $\rho=1/2$. Hence $LP_t^H$ is a feasible \textsc{LP-$\rho$-SEP} guide for $X_0$.
    \end{proof}
